% !TeX root = ../../Thesis.tex
\ThesisChapter
  {\ThesisText{Discussion}{Diskussion}}
  {ch:discussion}
  {TODO: Add a brief chapter overview. \\
  This chapter interprets the results, positions them against the related work, and examines their significance, limitations, and threats to validity.}

% TODO: Interpret the results presented in the evaluation chapter.

The evaluation chapter reports what was measured.
This chapter explains what those measurements mean.

Possible contents include:
\begin{itemize}
  \item the interpretation of each significant result, including the unexpected ones,
  \item the comparison of these results with the approaches surveyed in the related work chapter,
  \item an assessment of whether the objectives and requirements of the thesis were met,
  \item the practical implications of the findings, and
  \item the limitations of the proposed approach.
\end{itemize}

Explain similarities, differences, advantages, and the possible reasons for differing results.
Do the results confirm, extend, or contradict existing findings?

An unexpected or negative result is a legitimate outcome and is worth discussing carefully.
Explaining why something did not work is often more informative than reporting that something did.

Do not introduce new experimental results in this chapter.
New results belong in the evaluation chapter, where they are presented before being interpreted here.

\section{\ThesisText{Threats to Validity}{Validitätsbedrohungen}}
\label{sec:threats-to-validity}

% TODO: Discuss what could undermine the conclusions of this thesis.

A thesis that names its own weaknesses is more convincing than one that claims to have none.
An unacknowledged threat that an examiner spots is far more damaging than one that is stated openly.

Consider at least the following categories:

\begin{itemize}
  \item Internal validity: could the observed effect have a cause other than the one claimed?
  \item External validity: how far do the results generalize beyond the tested system, workloads, and scale?
  \item Construct validity: do the chosen metrics actually measure the property of interest?
\end{itemize}
Typical threats in systems and performance work include measurement noise with too few repetitions, results obtained on a single machine or cluster, a workload selection that happens to favor the proposed approach, and extrapolation beyond the scale that was actually measured.

State for each threat how it was mitigated, or why it could not be.
